\documentclass[a4paper,12pt]{article}
\usepackage[utf8]{inputenc}
\usepackage[russian]{babel}
\usepackage{amsmath, amssymb}


\begin{document}

\section*{Мультипликативные группы конечных полей}

Пусть $q = p^m$. Мультипликативная группа конечного поля $\mathbb{F}_q$ обозначается $\mathbb{F}_q^*$ и состоит из всех ненулевых элементов поля. Она является абелевой циклической группой порядка $q-1$.

\textbf{Пример.} В поле $\mathbb{F}_7$ множество $\{1,2,3,4,5,6\}$ образует группу порядка $6$. Пусть $\alpha = 3$. Тогда:
\[
3^1 = 3,\quad
3^2 = 9 \equiv 2,\quad
3^3 \equiv 6,\quad
3^4 \equiv 4,\quad
3^5 \equiv 5,\quad
3^6 \equiv 1.
\]
Таким образом, $\alpha$ является образующим элементом.

\section*{Порядок элемента}

Для элемента $a \in \mathbb{F}_q^*$ его порядок определяется как минимальное число $s > 0$, для которого
\[
a^s = 1.
\]
Из свойств циклических групп следует, что $s \mid (q-1)$.

\section*{Подгруппы}

Любая подгруппа $\mathbb{F}_q^*$ также циклическая, а её порядок делит $q-1$. Следовательно, порядок любого элемента является делителем порядка всей группы.

\section*{Аналог малой теоремы Ферма}

Для любого $x \in \mathbb{F}_q$ выполняется:
\[
x^q = x.
\]
Если $x \neq 0$, то это равносильно:
\[
x^{q-1} = 1.
\]

\section*{Лемма о порядке произведения}

Пусть $a,b \in \mathbb{F}_q^*$, причём $\operatorname{ord}(a)=s_a$, $\operatorname{ord}(b)=s_b$, и $\gcd(s_a, s_b)=1$. Тогда:
\[
\operatorname{ord}(ab) = s_a s_b.
\]

\textbf{Идея доказательства.} Сначала показывается, что $(ab)^{s_as_b}=1$, значит порядок делит $s_as_b$. Далее из равенства $(ab)^t=1$ выводится, что $t$ делится и на $s_a$, и на $s_b$, откуда следует утверждение.

\section*{Существование элементов заданного порядка}

Пусть число $s^t$ делит $q-1$, где $s$ — простое. Тогда существует элемент порядка $s^t$.

\textbf{Идея.} Если $\alpha$ — образующий элемент группы, то элемент
\[
\alpha^{\frac{q-1}{s^t}}
\]
имеет требуемый порядок.

\textbf{Пример.} В $\mathbb{F}_7$ имеем $q-1=6$. Возьмём делитель $2$. Тогда:
\[
\alpha = 3,\quad \alpha^{6/2} = 3^3 \equiv 6,
\]
и элемент $6$ имеет порядок $2$.

\section*{Порядки элементов в $\mathbb{F}_7^*$}

\[
\begin{array}{c|c|c}
a & \text{представление} & \operatorname{ord}(a) \\
\hline
1 & 3^0 & 1 \\
2 & 3^2 & 3 \\
3 & 3^1 & 6 \\
4 & 3^4 & 3 \\
5 & 3^5 & 6 \\
6 & 3^3 & 2 \\
\end{array}
\]

\section*{Примитивные элементы}

Существует элемент порядка $q-1$, который называется примитивным. Это напрямую следует из цикличности группы $\mathbb{F}_q^*$.

\section*{Коды максимальной длины}

Циклический код длины
\[
n = 2^m - 1,
\]
построенный с использованием примитивного полинома над $\mathbb{F}_2$, называется кодом максимальной длины.

Он имеет параметры $(2^m - 1, m)$ и является двойственным к коду Хэмминга. Порождающий многочлен можно представить в виде:
\[
g(x) = \frac{x^{2^m - 1} - 1}{h(x)},
\]
где $h(x)$ — примитивный полином.

Такие коды также известны как симплекс-коды.

\section*{Кодирование}

Для циклического $(n,k)$-кода кодовое слово получается как:
\[
c(x) = u(x)\,g(x),
\]
где $u(x)$ — информационный многочлен степени меньше $k$.

\textbf{Пример.} Код длины $15$ и размерности $4$ над $\mathbb{F}_2$ можно построить с использованием примитивного полинома:
\[
h(x) = 1 + x + x^4.
\]

\end{document}